% Companion Part VI -- VI/35 editorial bridge. Addition only; no existing proof replaced.
\paragraph{From homogeneous primes to affine charts.}
For a positively graded ring \(S\), the space
\(\operatorname{Proj}S\) consists of homogeneous prime ideals
\(\mathfrak p\) for which \(S_+\not\subseteq\mathfrak p\).
Its basic opens are
\[
 D_+(f)=\{\mathfrak p\in\operatorname{Proj}S:f\notin\mathfrak p\},
\]
where \(f\) is homogeneous of positive degree. The central local
computation is
\[
 \boxed{D_+(f)\cong\operatorname{Spec}(S_f)_0.}
\]
The \(\operatorname{Spec}\) on the right is an ordinary affine
scheme. Consequently, \(\operatorname{Proj}S\) is built by
\emph{gluing affine schemes}, not by inventing a different kind
of local ring.

\paragraph{Worked construction: projective line from two affine lines.}
Put \(S=k[x_0,x_1]\), with \(\deg x_i=1\). The chart
\(U_0=D_+(x_0)\) has coordinate
\(t=x_1/x_0\); the chart \(U_1=D_+(x_1)\) has coordinate
\(s=x_0/x_1\). Thus
\[
 U_0\cong\operatorname{Spec}k[t],\qquad
 U_1\cong\operatorname{Spec}k[s].
\]
On the overlap both \(x_0\) and \(x_1\) are invertible, and
\[
 k[t,t^{-1}]\cong k[s,s^{-1}],\qquad s=t^{-1}.
\]
This reciprocal coordinate change glues two affine lines into
\(\mathbb P^1_k\). Each chart looks affine, but the resulting
projective line is not affine: its global regular functions are only
\(k\), whereas an affine scheme with global ring \(k\) would
have a single point.

\paragraph{Worked conic: the grading leads to explicit chart equations.}
Let \(S=k[x,y,z]/(y^2-xz)\), with all variables of degree one.
The homogeneous relation defines a projective conic
\(C=\operatorname{Proj}S\). On the chart \(D_+(x)\), set
\(u=y/x\) and \(v=z/x\). Dividing the relation by \(x^2\)
gives \(u^2=v\); hence
\[
 (S_x)_0\cong k[u,v]/(u^2-v)\cong k[u].
\]
Similarly, on \(D_+(z)\) the coordinate \(w=y/z\) gives
\((S_z)_0\cong k[w]\). These two charts cover the conic because
if \(x=z=0\), the homogeneous equation forces \(y=0\), which
is excluded from projective space. On their overlap,
\[
 w=\frac{y}{z}=\frac{1}{u},\qquad
 u=\frac{y}{x},\quad z/x=u^2,
\]
so they glue by inversion. The projective conic is therefore
isomorphic to \(\mathbb P^1_k\). This concrete computation links
homogeneous equations, localizations, gluing, and projective geometry.


\paragraph{Why homogeneous equations give closed projective subschemes.}
Let \(I\subseteq S\) be a homogeneous ideal. The quotient
\(S/I\) is graded, and the projective subscheme cut out by \(I\)
is \(\operatorname{Proj}(S/I)\). For a homogeneous element \(f\),
its part inside \(D_+(f)\) is described by the ordinary affine
quotient of \((S_f)_0\) by the degree-zero elements of \(I_f\).
Thus a global homogeneous equation turns into familiar polynomial
equations after choosing projective coordinates.

For instance, in \(\mathbb P^1_k=\operatorname{Proj}k[x_0,x_1]\),
the ideal \((x_0)\) cuts out the point \([0:1]\).
On the chart \(x_1\ne 0\), set \(s=x_0/x_1\). The local equation
is \(s=0\), and its ring is \(k[s]/(s)\cong k\).
If we use the homogeneous ideal \((x_0^2)\) instead, the topological
support is still that same point, but the local equation is
\(s^2=0\). Consequently the local ring becomes
\[
 k[s]/(s^2),
\]
which has a nonzero nilpotent element \(s\). This \emph{double point}
is a projective scheme with a different infinitesimal structure,
even though its underlying set has not changed. The calculation
connects the graded quotient with the earlier discussion of
nonreduced schemes and fibers.

\paragraph{A projective point need not be a rational coordinate pair.}
When \(k\) is algebraically closed, closed points of
\(\mathbb P^1_k\) are commonly written \([a:b]\) with
\((a,b)\ne(0,0)\), modulo common nonzero scaling. Over a general
field, this familiar coordinate picture describes the
\(k\)-rational points, not all scheme-theoretic closed points.
For example, over \(\mathbb R\), the homogeneous polynomial
\(x_0^2+x_1^2\) generates a prime ideal of
\(\mathbb R[x_0,x_1]\); it determines a closed point of
\(\mathbb P^1_{\mathbb R}\) whose residue field is \(\mathbb C\).
On \(D_+(x_1)\), its coordinate ring is
\[
 \mathbb R[s]/(s^2+1)\cong\mathbb C.
\]
The scheme construction naturally includes this point even though
no real pair \([a:b]\) satisfies \(a^2+b^2=0\).

\paragraph{The structure sheaf remembers the graded input.}
On a basic open \(D_+(f)\), regular sections come from
\((S_f)_0\). At a point \(\mathfrak p\), the stalk is formed by
localizing homogeneous fractions of total degree zero at homogeneous
elements outside \(\mathfrak p\). This is the projective analogue
of the localization construction for \(\operatorname{Spec}A\):
local rings record functions near a point and make transition maps
compatible on chart overlaps.

In the following sections, projective space and projective schemes
will be studied through these same affine charts. The useful mental
picture is not that projective geometry discards affine algebra,
but that it uses affine algebra while keeping coordinate scale out
of the local descriptions.
